\cotefile{16}
\scanpage{44}
% archive sheet number: 43
\hfill\handwritten{4.1.67.}\\
Dear Coates,\\

I want to add a few more comments to the talk on algebraic cycles
and to what I told you on the phone.

I think the best will be to state the index conjecture right after
the statement of the main results of Hodge theory, adding that this conjecture
will take its whole significance only when coupled with ``conj.$A$ ''
in the next paragraph. This will give more freedom in the next paragraph
to express some extra relationships between various conjectures, such
as $A+\text{index}$ implies $B$.

In \deleted{c}\handwritten{\unsure{ch|c}}ar. zero, state some known extra features: index theorem holds,
the \deleted{conjectures}\inserted{\handwritten{\unsure{properties|propositions}}} $A_{\unsure{\ell|l|\chi}}$ to $D_{\ell}$ are indep\deleted{\unsure{nd}}endant of $\ell$ (because of the existence
of Betti cohomology, so that these \deleted{conjectures}\inserted{\handwritten{\unsure{properties|propositions}}} are equivalent to the
corresponding one's for rational cohomology), \deleted{and} $A$ and $C$ are independant
of the choosen polarization\deleted{\unsure{,|x}} (for $A$ because it is equivalent with $B$, for
$C$ because it can be expressed in terms of $A$, $C(X)=(A(X\mathbin{\mathrm{x}}X)+A(Y\mathbin{\mathrm{x}}Y)+\mathinner{..})$
Thus the \deleted{conjectures}\inserted{conditions} without ambiguity can be called $A(X)$ to $D(X)$,
without \deleted{index}\inserted{subscript} $\ell$ and without indication of polarization. Say too that
it is known that $C(X)$ is of finite dimension over \uline{$Q$}, \deleted{\unsure{we}}(so that $A$ can also
be expr\unsure{ë|e}ssed in terms of an equality of dimensions of $C^{i}$ and $C^{n-i}$\deleted{\unsure{)|$\chi$}} , which
again proves it is independant of \hspace{1em}), bu that this is not known in
car $p\handwritten{>}0$. Contrarily to what I hastily stated in my talk (influenced
from my recollections of the car $0$ case) it is not clear to me if in car
$p\handwritten{>}0$ the \deleted{\unsure{conjectures}}conditions \deleted{\unsure{$A(X$}} $A_{\ell}(X,\handwritten{\unsure{\xi|\zeta}})$ and $C_{\ell}(X,\handwritten{\unsure{\xi|\,\vert\,}})$ are independant of the
polarization $\handwritten{\xi}$; if you do not find some proof of this independance, then
the possible dependance should be pointed out, as well as the fact that \deleted{a}
\deleted{priori} we do not have a proof that $A$ to $D$ are independant of $\ell$. Of cour%
\scanpage{45}%
% archive sheet number: 44
se, if the index theorem is proved for $X$, then $\deleted{\unsure{D|O}}A_{\unsure{\ell|l|\chi}}(X,\handwritten{\unsure{\xi|\zeta}})=B_{\ell}(X)$ is again
independant of the polarization, and analoguous remark for $C_{\ell}(X,\handwritten{\xi})$.

When speaking about condition $C_{\ell}(X,\handwritten{\xi})$, emphasize at once its stability
properties by products (the proof I suggested works indeed) specialization
(with possib\deleted{\unsure{h}}le change of characteristics), hyperplane or more generally
linear sections. Give an extra proposition for the relations with the
property $A$, via a formal proposition as follows:

\uline{Proposition}\quad Condititions équivalentes sur $X$(varieté \uline{polarisée}):

\[
C_{\ell}(X) \tag*{(i)}
\]

\[
C_{\ell}(Y)\text{ et }A_{\ell}(X\mathbin{\mathrm{x}}X)^{\deleted{\unsure{o|0}}},\quad
\deleted{\text{où l'exposant }{}^{o}} \tag*{(ii)}
\]

(ii bis)\quad $C_{\ell}(Y)$ \deleted{\unsure{e}}et $A_{\ell}(X\mathbin{\mathrm{x}}X)^{\unsure{o|0}}$, où l'exposant $^{o}$ signifie qu'on se borne
à exprimer la condition $A$ pour l'homomorphisme en dimension c\deleted{\unsure{ir}}ritique
$H^{2n-2}\handwritten{\longrightarrow}H^{2n+2}$.

\[
C_{\ell}(Y)\text{ et }A_{\ell}(X\mathbin{\mathrm{x}}Y). \tag*{(iii)}
\]

(iii bis)\quad $C_{\ell}(Y)$ et $A_{\ell}(X\mathbin{\mathrm{x}}Y)^{o}$, où l\deleted{\unsure{b}}'exposant $^{o}$ signifie qu'on se borne
à exprimer la condition $A$ en dimension c\deleted{\unsure{ir}}ritique $H^{(2n-1)-1}\handwritten{\longrightarrow}H^{(2n-1)+1}$.

(iv)\quad \deleted{\unsure{Pour tout $0$|Pour tout $O$}} $C_{\ell}(Y)$, et pour tout $i\handwritten{\leq} n-1$, l'homomorphisme
naturel $H^{i}(Y)\handwritten{\unsure{-|\to}}H^{i}(X)$ inverse à gauche de $\phi^{i}:H^{i}(X)\handwritten{-}H^{i}(Y)$ (induit
par $\handwritten{\Lambda}_{X}\phi_{\handwritten{*}}$) est induit par une classe de correspondance algébrique (induisant ce qu'elle veut sur les autres $H^{j}(Y)$ ).

(iv bis)\quad $C_{\ell}(Y)$, et pour $j\handwritten{\geq} n+1$, l'homomorphisme naturel $H^{j}(X)\handwritten{\to}H^{\unsure{j-2}}(Y$
inverse à droite de $\phi_{j-2}:H^{j-2}(Y)\handwritten{\unsure{-|\to}}H^{j}(X)$ \deleted{est} (induit par $\phi^{*}\handwritten{\Lambda}_{X}$) est
induit par une classe de correspondance algébrique (induisant ce qu'elle
veut sur les autres $H^{i}(X)$).

\uline{Corollaire}\quad \deleted{\unsure{L}}Ces conditions équivalent aussi à

(v)\quad $A_{\ell}(X\mathbin{\mathrm{x}}X)^{\unsure{o|\,|0}}+A_{\handwritten{\ell}}(Y\mathbin{\mathrm{x}}Y)^{\handwritten{o}}+A_{\handwritten{\ell}}(Z\mathbin{\mathrm{x}}Z)^{\handwritten{o}}+\hspace{2em}$, où $X\handwritten{\supset}Y\handwritten{\supset}Z$ est une
suite décorissante de sections hyperplanes\handwritten{\unsure{ non sing.}}

\scanpage{46}
% archive sheet number: 45
(vi)\quad $A_{\unsure{\ell|l|\chi}}(X\mathbin{\mathrm{x}}Y)^{\unsure{o|0}}+A_{\ell}(Y\mathbin{\mathrm{x}}Z)^{o}+\mathinner{....}$, avec les mêmes notations.

Of course, the products and hyperplane sections are endowed with the
polarizations stemming from the polarization on $X$. The conditions (v) and
(vi) have t\deleted{\unsure{h}}the slight interest that they allow to express the conjecture
$A(k)=C(k)$ in te\deleted{\unsure{r}}rms of $A(T)^{o}$ for every $T$ of even (resp. odd dimension),
where the upper $^{o}$ means that it is sufficient to look at what happens in
critical dimensions.

For thenproof of the proposition, I told you already the equivalence
of (i) and (ii),(ii bis). The equivalence of (iv) and (iv bis) is
trivial by transposition, they \deleted{are of course implied by (i) because of
the stabilities, because $C_{\ell}(X)$ \unsure{$\Rightarrow$} $C_{\ell}(Y)$,} imply (i) because
$H^{2n-i}(X)\handwritten{\to}H^{i}(X)$ is the composit\deleted{\unsure{i}}um
$H^{2n-i}(X)\handwritten{\unsure{-|\to}}H^{2n-i-2}(Y)\handwritten{-}H^{i}(Y)\handwritten{-}H^{i}(\unsure{X}$
where the extreme arrows are the ones of (iv bis) and (iv) and the
middle one is induced by $\handwritten{\Lambda}_{Y}^{(n-1)-i}$, and they are implied by (iii bis)
be\deleted{\unsure{e}}cause of the formula
\[
\deleted{\unsure{\Lambda_{X}\phi\phi}}\quad
(\handwritten{\Lambda}_{X}\phi_{\handwritten{*}})L_{Y}+L_{X}(\handwritten{\Lambda}_{X}\phi_{\handwritten{*}})
=(\phi_{\handwritten{*}}\handwritten{\Lambda}_{Y}\phi^{\handwritten{*}}+\mathrm{id}_{X})\phi_{\handwritten{*}}.
\]
\inserted{\handwritten{N.B. \((\varphi^{*}\Lambda_{X})L_{X}+L_{Y}(\varphi^{*}\Lambda_{X})=\varphi^{*}(\varphi_{*}\Lambda_{Y}\varphi^{*}+\unsure{\mathrm{id}_{X}})\)}}
On the other hand (iii)$\handwritten{\Longrightarrow}$(iii bis) is trivial, and so is (i)$\handwritten{\Longrightarrow}$(iii)
because of the stabilities.

For the list of the known facts, you can state that :

1) In arbitrary c\deleted{\unsure{a}}aracteristics, $C(X)$ is known if $\dim X\handwritten{\leq}2$, because
more generally, it is known that in arbitray dimension $n$, $H^{2n-1}(X)\handwritten{\to}H^{1}(X)$
is induced by an algebraic correspondance class; also, in arbitray dimension,
it is known that $\handwritten{\pi}_{o},\handwritten{\pi}_{2n},\handwritten{\pi}_{1},\handwritten{\pi}_{2n-1}$ are algebraic (trivial for
the first two, not quite trivial for the two next one's). If $\dim X=3$,
it is not known however, \deleted{\unsure{ve|ev}}even in \deleted{\unsure{dim}}\unsure{car|dim} $0$, if $C(X)$ or only $D(X)$ hold,
nor $A(X)$ and $B(X)$ in car. $p\handwritten{>}0$\handwritten{; also if for $1$-cycles, $\tau$-equ. $=$
num. \unsure{équivalence|equivalence} ....}

\scanpage{47}
% archive sheet number: 46
By the way, the fact that the $\handwritten{\pi}_{i}$ for a surface are algebraic was
pointed out (Tate tells me) by Hodge in Algebraic correspondances between
surfaces, Proc. London Math. So\deleted{\unsure{6}}c. Seeies 2, Vol XLIV, 1938, p.226. It is
rather striking that this statement should not have struck the algebraic
geometers more, and has fallen into oblivion for \inserted{\handwritten{\unsure{nearly}}} thirty years !

2) In car.$0$, $A(X)$ is known for $\dim X\handwritten{\leq}4$. But $A(X)^{\unsure{o|0}}$ is not known if
$\dim X=5$; the first interesting case would be for a variety $X\mathbin{\mathrm{x}}Y$, $X$
of $\dim 3$ and $Y$ a hyperplane section, as this would prove $C(X)$, see above.

Thus the main problems arise already for $1$-cycles on threefolds,
and partially even in car.$0$. Urged by Kleiman\deleted{n}'s question, I wi\deleted{\unsure{h}}ll look
again at my old scribbles on that subject (when I pretend to reduce the
``strong'' form of Lefschetz to the ``week'' one). As for the suggestion I
made on the phone, to try to get any $X$ as birat\deleted{\unsure{n}}ionally equivalent to a
non singular $X'$, which is a spcialisation of a non singular $X''$, itself
birationally equivalent to a non singular hypersurface - this cannot
work as Serre pointed out, because such an $X$ would have to be simply connected !
Thus if one wants to reduce someh\deleted{\unsure{eu}}ow to the case of hypersurfaces,
one will have to work also with singular ones, and see how to reformulate
for singular varieties the standard conjectures ...

Sincerely your's\\

\scanpage{48}
% archive sheet number: 47
\hfill\handwritten{6.1.1966}\\
Dear Coates,\\

Here a\handwritten{$\vert$}few more comments to my talk on the conjectures. The following
proposition shows that the conjecture $C_{\unsure{\ell|l|\chi}}(X)$ is independant of the choosen
polarization, and has also some extra interest, in showing the part playd
by the fact that $H^{i}(X)$ should be ``motive-theor\deleted{\unsure{i}}etically'' isomorphic to
its natural dual $H^{2n-i}(X)$ (as usual, I drop the twist for simplicity).

\uline{Proposition}\quad The condition $C_{\unsure{1|l}}(X)$ is equivalent also to \inserted{each of} the follwing
conditions:

a)\quad $D_{\ell}(X)$ ho\deleted{r}lds, and for every $i\inserted{\handwritten{<}n}$, there exists an isomorphism
$H^{2n-i}(X)\handwritten{\to}H^{i}(X)$ which is algebraic (i.e. induced by an algebraic correspondance
class; we do not make any assertion on what it induces in degrees
different from $2n-i$).

b)\quad For every endomorphism $H^{i}(X)\handwritten{\unsure{-|\to}}H^{i}(X)$ which is algebraic, the
coefficient\unsure{s} of the characteristic polynomial are rational, and for every
$i\inserted{\handwritten{<}n}$, there exists an isomorphism $H^{2n-i}(X)\handwritten{-}H^{i}(X)$ which is algebraic.

Proof. I sketched already how $D_{\ell}(X)$ implies the fact that for an
algebraic endomorphism of $H^{i}(X)$, the coefficients of the characteristic
polynomial are rational numbers, Therefore we know that a) implies b),
and of course $C_{\ell}(X)$ implies a). It remains to prove that b) implies
$C_{\ell}(X)$. Let $u:H^{2n-i}(X)\handwritten{-}H^{i}(X)$ be the given isomorphism which is algebraic,
and $v:H^{i}(X)\handwritten{-}H^{\deleted{\unsure{a}}2n-i}(X)$ \deleted{an}\handwritten{the} algebraic isomorphism in the opposite
direction, induced by $L_{X}^{n-i}$. Then $uv=w$ is an automorphism of $H^{i}(X)$ which
is algebraic, a\deleted{\unsure{ad}}nd the Hamilton-Cayley formula $\unsure{u|w}^{\unsure{h|b}}-\handwritten{\sigma}_{1}(w)u^{h-1}+\mathinner{...}+\deleted{(-1)}$
$(-1)^{b}\handwritten{\sigma}_{b}(w)=0$ (where the $\handwritten{\sigma}_{i}(w)$ are the coefficients of the car. pol.
of $w$) whos that $w^{-1}$ is a linear combination of the $w^{i}$, with coe\deleted{\unsure{e}}fficient
of the type $\pm\handwritten{\sigma}_{i}(w)/\handwritten{\sigma}_{b}(w)$ (NB $b=\operatorname{rank}H^{i}$). The assumption implies
\scanpage{49}
% archive sheet number: 48
that these coefficients are rational, which implies that $w^{-1}$ is algebraic,
and so is $w^{-1}u=v^{-1}$, which was to be proved.

NB\quad In char. $0$, \deleted{it is} the statement simplifies to: $C(X)$ equivalent
to the existence of algebraic isomorphisms $H^{2n-i}(X)\handwritten{\to}H^{i}(X)$,(as the preliminar\deleted{\unsure{a}}y
condition in b) is then automatically satisfied). Maybe with some
extra care this can be proved too in arbitrary characteristics.

\uline{Corollary}\quad Assume $X$ and $X'$ satisfy condition $C_{\unsure{\ell|l|\chi}}$, and let
$u:H^{i}(X)\handwritten{\unsure{-|\to}}H^{i+2\deleted{p}j}(X')$ \inserted{$(\deleted{p}j\handwritten{\in}\text{\uline{$Z$}})$} be an isomorphism which is algebraic. Then $u^{-1}$ is
algebraic.

Indeed, the two spaces can be identified ``algebraically'' (both directions !) to their dual, so taht the transpose of $u$ can be viewed as an
isomorphism $u':H^{i+2\deleted{p}j}(X')\handwritten{-}H^{i}(X)$. Thus $u'u$ is an \inserted{algebraic} automorphism \inserted{$w$} of
$H^{i}(X)$, and by the previous argument we see that $w^{-1}$ is algebraic, hence
so is $u^{-1}=w^{-1}u'$.

As a consequence, w\unsure{ë|e} see that if $x\handwritten{\in}H^{i}(X)$ is such that $u(x)$ is
algebraic ($i$ being now assumed to be even), then so is $x$. The same result
should hold in fact if $u$ is a monomorphism, the reason being that in this
case there should exist a left-inverse which is algebraic; this exists indeed in a case like $H^{n-1}(X)\handwritten{\to}H^{n-1}(Y)$ (where we take the left inverse
$\handwritten{\Lambda}_{X}\phi_{\handwritten{*}}$). But to get i\deleted{\unsure{n}}t in general, it seems we need moreover the
Hodge index relation. (The complete yoga then being that we have the
category of motives which is semi-simple !). Without speeking of motives,
and staying down on earth, it would be nice to explain in the notes that
$C(X)$ together with the index relation $I(X\mathbin{\mathrm{x}}X)$ implies that the ring of
correspondance classes for $X$ is semi-simple, and how one deduces from this
the existence \deleted{\unsure{an}}of left and right inverses as looked for above.

\scanpage{50}
% archive sheet number: 49
This could be given in an extra paragraph (which I did not really touch
upon in the talk), containing also the deduction of the Weil conjectures
from the conjectures $C$ and $A$.

A last and rather trivial remark is the following. Let's introduce
\deleted{\unsure{a|$x$}} variants $A'_{\unsure{\ell|l|\chi}}(X)$ and $A''_{\ell}(X)$ as follows:

$A'_{\ell}(X)$ : if $2i\handwritten{\leq} n-1$, any element \inserted{$x$} of $H^{i}(X)$ whose image in $H^{i}(Y)$ is alge\handwritten{\unsure{(|\,}}braic, is algebraic.

$A''_{\ell}(X)$ : if $2\deleted{i}\unsure{j|i}\handwritten{\geq} n-1$, any \inserted{algebraic} element \inserted{\deleted{\unsure{$x$}}}\deleted{\unsure{of} $H^{\unsure{i|j}}(Y)$ whose image \unsure{in|on}}\handwritten{of}
$H^{\deleted{i}j+2}(X)$ is \deleted{algebraic, is algebraic.}\deleted{\handwritten{\unsure{is}}} the image of an algebraic element of $H^{i}(Y)$.

Let us consider also the specifications $A'_{\ell}(X)^{\unsure{o|0}}$ and $A''_{\ell}(X)^{o}$, where we
restrict to the \unsure{crittval|critival|critiwal} dimensions $2i=n-1$ if $n$ odd, $2i=n-2$ if $n$
even. All these conditions are in the nature of ``week'' Lefschets relations,
and they are trivially implied by $A_{\ell}(X)$ resp. $C_{\ell}(Y)$ (in the first
case, applying $\phi$ we see that $L_{X}x$ is algebraic; in the second, we take
$y=\handwritten{\Lambda}_{Y}\phi^{\handwritten{*}}(x)$). The remark then is that these pretend\inserted{\handwritten{ed}}ly ``week'' variants
in fact imply the full Lefs\deleted{\unsure{b}}chets relations for al\inserted{\handwritten{ge}}braic cycles, namely:

\uline{Proposition}\quad $\deleted{\unsure{A}}\unsure{C|E}_{\ell}(X)$ is equivalent to the conj\deleted{\unsure{an}}unction $C_{\ell}(Y)+A'_{\ell}(X\mathbin{\mathrm{x}}X)^{o}+{}$
$A''_{\ell}(X\mathbin{\mathrm{x}}X)^{o}$, hence (by induction) also to the conj\deleted{\unsure{an}}unction of the
conditions $ {A'_{\ell}}^{o}$ and $ {A''_{\ell}}^{o}$ for all of the varieties $X\mathbin{\mathrm{x}}X,Y\mathbin{\mathrm{x}}Y,Z\mathbin{\mathrm{x}}Z$ .... . Analoguous statement with $X\mathbin{\mathrm{x}}Y,Y\mathbin{\mathrm{x}}Z$ etc instead of \uline{$X\mathbin{\mathrm{x}}X,Y\mathbin{\mathrm{x}}Y$ etc.}

This comes from the remark that \deleted{if $X$ is even dimensional (say a
square like \unsure{$X\mathbin{\mathrm{x}}X$}) then} $A_{\ell}(X\deleted{\unsure{\mathbin{\mathrm{x}}X}})^{o}$ follows from the conjuction of
$A'_{\ell}(X)^{o}$ and $A''_{\ell}(X)^{o}$, as one sees by decomposing $L_{X}^{2}:H^{2m-2}(X)\handwritten{\to}H^{2m+2}(X)$
into $H^{2m-2}(X)\handwritten{\xrightarrow{\varphi^{*}}}H^{2m-2}(Y)\handwritten{\xrightarrow{\varphi_{*}}}H^{2m}(X)\handwritten{\overset{L_{X}}{-}}H^{2m+2}\deleted{\unsure{\mathrm{XX}}}(X)$ if $\dim X=2m$ is even,
and $H^{2m\deleted{\unsure{(}}+1-1}(X)\handwritten{\to}H^{2m+1+1}(X)$ into $H^{2m}(X)\handwritten{\overset{\varphi^{*}}{-}}H^{2m}(Y)\handwritten{\overset{\varphi_{*}}{\unsure{-|\to}}}H^{2m+2}(X)$ if
$\dim X=2m+1$ is odd.

Sincerely your's\\

\scanpage{51}
% archive sheet number: 50
\uline{Lemme}\quad Soit $\unsure{\mathcal{A}|A}$ un anneau d'endomorphismes
du \unsure{groupe|grpe|gpe} \deleted{\unsure{a}}abélien $H^{*}(X)$, stable par composantes homogènes
(degré \uuline{total}), par transposition. \deleted{\unsure{et tel}} \inserted{Supposons $\unsure{\xi|\zeta}\in\mathcal{A}$}, et que pour
tout $0\leq i<n$, il existe $u_i\in\mathcal{A}$, \inserted{de degré $\deleted{\unsure{2n}}-2(n-i)$,} tel que
$u_i\xi^{n-i}:H^i(X)\longrightarrow H^{2n-i}(X)\longrightarrow H^{\deleted{\unsure{2}}i}(X)$ soit l'identité.
Alors $\mathcal{A}$ contient $L$\margin{(\unsure{et réciproquement|et réciproquemt|et vice versa})}, \deleted{\unsure{les et}} et les projections
% slanted
sur les facteurs \unsure{élémentaires} $\xi^j P^i$ ($i\leq n$, et
\deleted{\unsure{et} $i+2j\leq n$ \unsure{ou} $i+2j\leq 2n-i$ \unsure{i.e.}} $j\leq n-i$)

\uline{Corollaire}\quad Conditions équivalentes :

(i) $L$ algébrique

(ii) pour \unsure{tt|et} $1\leq i\unsure{<|\leq}n$, il existe une corr. \uline{alg}
induisant $H^{2n-i}(X)\unsure{-|\to}H^i(X)$ inverse de $\unsure{J|f|\zeta}^i$

\deleted{\unsure{(iii)}} S'il en est ainsi, alors les \deleted{\unsure{$\xi$}} projections
sur les facteurs \unsure{élémentaires} de $H^{*}(X)$
sont algébriques, et \unsure{les} projections $\deleted{\unsure{de}}H^{*}(X)\unsure{\to|-}H^i(X)$
aussi.

\uline{\unsure{Dém|Dém.} du lemme.}\quad Comme $\mathcal{A}$ stable par
transposition, il suffit de regarder les $\unsure{\xi^jP^i|\zeta^jP^i}$,
pour $i+2j\leq n$. Soit $k=i+2j$, on procède
par récurrence sur $k$, trivial si $k<0$.

\scanpage{52}
% archive sheet number: 51
Algèbres étales et la différente \mspage{(4)}\hfill Tribu n\textsuperscript{o} 65

\S~2. \uline{Différente}.

\uline{n\textsuperscript{o} 1}. \uline{Idéaux de Filting d'un module} $M$.

Voir thèse de D. Lazard. Lien avec le contenu. Le premier Filting est compris entre
$\operatorname{Ann}(M)$ et une de ses puissances.

\uline{n\textsuperscript{o} 2}. \uline{Différente de Kähler}.

On prend pour $B$ une algèbre de p. f. sur $A$ (plus généralement il suffit peut
être de supposer que $\handwritten{\Omega}_{A}(B)$ est un $B$-module de p. f.). Par définition la différente
de Kähler $\text{\uline{$D$}}_{B/A}$ est le premier idéal de Filting de $\handwritten{\Omega}_{A}(B)$ :
\[
\text{si }\quad A^{p}\xrightarrow{u}A^{q}\longrightarrow\handwritten{\Omega}_{A}(B)\longrightarrow 0,\quad \text{c'est }\operatorname{Im}(\handwritten{\unsure{\Lambda|\bigwedge}}^{q}u).
\]
La variété des zéros de $\text{\uline{$D$}}_{B/A}$ est l'ensemble des points de ramification. Comportement
de $\text{\uline{$D$}}_{B/A}$ par changement d'anneau (c'est trivial).

\uline{n\textsuperscript{o} 3}. \uline{Transitivité}.

Dire ce qu'on sait. Si on a $A\to B\to C$ avec $C$ intersection complète locale
sur $B$, la formule de transitivité marche.

\uline{n\textsuperscript{o} 4}. \uline{Lien avec la trace}.

Définir la différente de Dedekind par la trace, et montrer qu'elle coïncide avec
celle de Kähler quand $A$ et $B$ sont des anneaux de valuation discrète (il s'agit
du gros th. 3 du n\textsuperscript{o} 7, \S~2, avec son dévissage, - éventuellement agrémenté des lemmes
du n\textsuperscript{o} 6 s'ils n'ont pas été faits dans un chapitre antérieur pour illustrer la notion
d'intersection complète. En tous cas on aura dit que si $A$ et $B$ sont réguliers,
$A\longrightarrow B$ est une intersection complète relative. La prop. 11 sera éclairée par le lemme
suivant :

``Si $k\longrightarrow R$ est finie et si $R/\operatorname{rad}(R)$ est étale, alors $R$ monogène ssi les
idéaux maximaux de $R$ sont homogènes''.

La prop. 12 est absorbée par le résultat ci-dessus sur les anneaux réguliers.

% Unlabelled faint outline in the left margin.
\figure{\begin{tikzpicture}[x=1mm,y=1mm]
\draw[gray,thin] (2.3,4.2) .. controls (1.9,4.9) and (1.1,4.5) .. (1.3,4.0)
 .. controls (1.5,3.7) and (1.9,3.8) .. (2.0,4.0)
 .. controls (1.5,3.7) and (.7,3.8) .. (.4,4.0)
 .. controls (.1,4.8) and (.8,5.8) .. (1.6,5.8)
 .. controls (2.2,5.8) and (2.3,5.1) .. (2.8,5.7)
 .. controls (3.4,6.6) and (3.5,6.5) .. (4.0,6.0)
 .. controls (4.4,5.7) and (4.7,6.4) .. (4.9,6.6)
 .. controls (5.1,6.8) and (5.1,6.2) .. (5.5,6.3)
 .. controls (5.8,6.8) and (6.1,5.8) .. (6.2,5.5)
 .. controls (6.4,5.1) and (6.2,4.7) .. (5.9,4.4)
 .. controls (6.7,4.6) and (7.2,4.4) .. (7.3,3.7)
 .. controls (7.5,3.9) and (7.8,3.3) .. (7.9,3.1)
 .. controls (8.4,2.4) and (8.2,1.0) .. (8.7,.4);
\end{tikzpicture}}

\uline{n\textsuperscript{o} 5}. \uline{Compléments}.

La ramification modérée du n\textsuperscript{o} 5 (noter que la formule de réduction des traces
utilisée en haut de la p. 25 est faite en Alg. VIII, \S~11).

Donner la description Kummerienne des extensions totalement et modérément ramifiées
(extraire la racine $n$-ième d'une uniformisante bien choisie ; ceci dans le cas complet);
exemple des surfaces de Riemann (ici le groupe de Galois a un générateur canonique
$\exp\left(\dfrac{2i\handwritten{\pi}}{n}\right)$ ).

\scanpage{53}
% archive sheet number: 52
Supposons $k\geq 0$, et \unsure{la|le} \unsure{dém|lemme} \unsure{fournie|prouvé} pour $k'<k$.

On sait déjà que les $p_{\deleted{\unsure{a}}i}:H^{\unsure{*|\ell}}(X)\unsure{-|\to}H^{\unsure{i|\ell}}(X)$
sont $\in\unsure{\mathcal{A}\mathrm{lg}|\mathrm{Alg}}$ pour $i<k$, donc (transp.) pour
$i>2n-k$. Soit $T_k=\mathrm{id}-\sum_{i<k}p_i-\sum_{j>2n-k}p_j$, donc
$T_k\in\unsure{\mathcal{A}|A}$, plus \unsure{gén.} $T_{k'}\in\mathcal{A}$ si $k'\unsure{\leq|<}k$. Alors
$\unsure{\xi|\zeta}^{n-k}T_k$ est nul sur $H^{\unsure{i|j}}$, $j\neq k$, et induit
$H^k\xrightarrow{\overset{\xi^{n-k}}{\sim}}H^{2n-k}$, donc $u_{\unsure{k|n}}\xi^{n-k}T_k=p_k$.

D'autre part, si $i+2j=k$\inserted{, $j>0$}, la projection sur
$\xi^j\unsure{P|p}^{i}$ est aussi $\in\mathcal{A}$, comme on voit \deleted{\unsure{ainsi|aussitôt}}
par récurrence ascendante sur $i$ (donc descendante \unsure{sur|en}
$j$) ainsi : si $i<0$ rien à prouver. Si c'est
fait pour les $i'<i$, on forme
$\xi^{n-k+j}T_k$, \deleted{\unsure{$H^j$}} qui induit \unsure{zéro|0|nul} sur $H^\alpha$ pour $\alpha\neq k$,
et induit $H^{k}\deleted{\unsure{\xi}}\xrightarrow{\xi^{\unsure{n-k-j|n-k+j}}}H^{2n-k+2j}$, \deleted{\unsure{qui|facteurs}}\inserted{qui est nul}
sur les facteurs $\xi^{j'}P^{i'}$ pour $j'<j$, \inserted{de somme $N$,} et sur
les facteurs $\xi^{j'}P^{i'}$ avec $j'\geq j$ (i.e. $i'\leq i$),
\deleted{\unsure{coïncide avec} $1$} dont la somme n'est autre
que $\xi^j H^{k-2j}$, son composé avec l'\unsure{isom}
$H^{k-2j}\xrightarrow{\xi^{j}}\xi^j H^{k-2j}$ est $\xi^{\unsure{k|n}-(k-2j)}$. Donc l'\unsure{hom|homom.}
$u_{k-2j}\xi^{n-k+j}T_k$\unsure{:|,|;}, nul sur $N$, est sur $\xi^j H^{k-2j}$
l'inverse de $H^{k-2j}\xrightarrow{\sim}\xi^j H^{k-2j}$. Donc

\scanpage{54}
% archive sheet number: 53
Algèbres étales et la différente \mspage{(3)}\hfill Tribu n\textsuperscript{o} 65

sur $A$ tel que, en posant $B'=A[x]$, $F'(x)$ soit inversible dans $B'_{\handwritten{\unsure{q|\mathfrak{q}}\cap}B'}$ et
$B'_{\handwritten{q\cap}B'}\longrightarrow B_{\handwritten{q}}$ soit bijectif. Généralisation au cas où $B$ est localisée et une
$A$-algèbre finie.

Si, de plus, $A_{\handwritten{q\cap}A}$ est normal et si $B$ est étale en $\handwritten{q}$, $B_{\handwritten{q}}$ est normal.
La prop. 4 (p. 4) du texte.

Ajouter, entre astérisques, les améliorations dépendant du main-theorem.

\uline{n\textsuperscript{o} 5}. \uline{Groupes d'automorphismes}.

La prop. 6 du texte, dans le cas noethérien non réduit (cf. Tribu n\textsuperscript{o} 64). Son cor.
(p. 7) est la prop. appliquée à $H$, de sorte qu'on le vide.

\uline{n\textsuperscript{o} 6}. \uline{Un critère de non ramification}.

Il s'agit de l'actuel n\textsuperscript{o} 4 (\S~1) amélioré comme suit. Remplacer (ici et partout)
l'idiote formule $(A^{*})^{-1}B=L$ par $KB=L$. Référer à Alg. IX, \S~2 pour la définition
du discriminant.

Dans le b) du th. 1, ajouter

(iii)\qquad $B$ est finie sur $A$ et $B/\handwritten{\unsure{m|\mathfrak{m}}}B$ est étale sur $A/\handwritten{m}$ ;

(iv)\qquad $B$ est étale finie sur $A$ ;

comme conditions équivalentes à (i) et (ii). Ça absorbe la partie intéressante de c),
ainsi que le cor. 1.

Pour démontrer le lemme (p. 10-11) on remarque que $P(x)=D(1,x,\ldots,x^{n-1})$ est
fonction polynôme de $x$, et qu'on a $P\handwritten{\not}=0$ (passer à la clôture algébrique). On a
donc $x\handwritten{\in}R$ tel que $P(x)\handwritten{\not}=0$.

Le cor. 1 est absorbé. La déf. 2 est supprimée. Le cor. 2 devient une prop. donnant
l'équivalence de : a)\quad $B$ est étale finie sur $A$ ;

b)\quad $B_{\handwritten{m}}$ est étale finie sur $A_{\handwritten{m}}$ pour tout idéal maximal $\handwritten{m}$ de $A$ ;

c)\quad $B$ est finie sur $A$ et $B/\handwritten{m}B$ est étale sur $A/\handwritten{m}$ pour tout
idéal maximal $\handwritten{m}$ de $A$.

Le cor. 3 est vidé. On donne enfin le cor. 4 (y remplacer ``non-ramifié'' par
``étale'').

On essaiera de mettre l'\uline{appendice} suivant : si $I$ est un idéal nilpotent de $A$,
il y a une équivalence de catégories entre les algèbres étales finies sur $A$ et les
algèbres étales finies sur $A/I$. Voir S. G. A.

\scanpage{55}
% archive sheet number: 54
$\unsure{\xi|\zeta}^{j}u_{k-2j}\xi^{n-k\unsure{-|+}j}T_k$ n'est autre que la
projection de $H^{\unsure{*|k|\alpha}}$ sur $\xi^j H^{k-2j}$, somme des
$\xi^{j'}\unsure{P|p}^{i'}$ pour $i'+2j'=k$, $j'\geq j$ ; \deleted{\unsure{or|car}} comme
la projection \deleted{\unsure{sur}} \unsure{partielle|partielles} sur $\xi^{j'}P^{\unsure{i|i'}}$,
$j'>j$, est $\in\unsure{\mathcal{A}|A}$, il en est de $=$ de
la projection sur $\xi^j P^i$, par différence.

Cela \unsure{prouve} le lemme.

\scanpage{56}
% archive sheet number: 55
Algèbres étales et la différente \mspage{(2)}\hfill Tribu n\textsuperscript{o} 65

\uline{n\textsuperscript{o} 2}. \uline{Algèbres étales et algèbres non ramifiées}.

Soit $B$ une $A$-algèbre. On considère les propriétés suivantes :

a)\quad $B$ est de p. f.

b)\quad $B$ est plate

c) Pour tout $\handwritten{\unsure{q|\mathfrak{q}}\in}\operatorname{Spec} B$ et pour $\handwritten{\unsure{p|\mathfrak{p}}}= \handwritten{q\cap}A$, $B_{\handwritten{q}}/\handwritten{p}B_{\handwritten{q}}$ est une extension
séparable finie de $A_{\handwritten{p}}/\handwritten{p}A_{\handwritten{p}}$.

En présence de a) \deleted{et b)}, c) veut dire que $\handwritten{\Omega}_{B/A}$ est nul en $\handwritten{q}$.

On dit que $B$ est étale si on a a), b), c).

On dit que $B$ est étale en $\handwritten{q\in}\operatorname{Spec}(B)$ s'il existe $f\handwritten{\notin q}$ tel que on ait a),
b), c) pour $B_f$. (i. e. tel que $B_f$ soit étale).

On dit que $B$ est non ramifiée si elle est de p. f. et si $\handwritten{\Omega}_{B/A}=0$.

Ainsi : étale $=$ non ramifiée $+$ plate.

Si $A$ est réduite, $B$ est réduite.

Transitivité ($A\longrightarrow B\longrightarrow C$) de l'étalité et de la non-ramification.

Effet d'un changement d'anneau sur ces notions : la montée marche dans le cas
général, la descente marche dans le cas fidèlement plat.

\uline{n\textsuperscript{o} 3}. \uline{Algèbres finies étales et non ramifiées}.

Les équivalences de la prop. 1, p. 1, du texte (améliorées comme dans la Tribu
n\textsuperscript{o} 64 ; on supposera seulement $B$ de présentation finie sur $A$ ).

Soient $B$ de p. f. sur $A$ noethérien, $\handwritten{q\in}\operatorname{Spec} B$ et $\handwritten{p}=\handwritten{q\cap} A$ ; pour
que $B$ soit étale en $\handwritten{q}$, il faut et il suffit que $\handwritten{\widehat}{B}_{\handwritten{q}}$ soit étale finie sur
$\handwritten{\widehat}{A}_{\handwritten{p}}$ ; analogue non ramifié ; la démonstration utilise la platitude supérieure. Corollaire pour le cas où les corps résiduels sont égaux (cf. S. G. A.).

Pour que $B$ soit étale finie sur $A$ connexe, il faut et il suffit qu'il existe
$A'$ fidèlement plate sur $A$ telle que $A'\handwritten{\otimes}B$ soit un produit de copies de $A'$
(si $A$ n'est pas connexe, ça ne marche que par morceaux).

La prop. 5 du texte (arrangée comme dans la Tribu n\textsuperscript{o} 64). Le calcul de traces de
la prop. 3, a) du \S~2 (p. 16).

\uline{n\textsuperscript{o} 4}. \uline{Structure locale des algèbres non ramifiées et étales}.

Soit $B$ une $A$-algèbre finie et soit $\handwritten{q\in}\operatorname{Spec}(B)$. Pour que $B$ soit non
ramifiée en $\handwritten{q}$ il faut et il suffit qu'il existe $x\handwritten{\in} B'$ de polynôme minimal $F$
